% 乔治·基斯佳科夫斯基（George Kistiakowsky）（综述）
% license CCBYSA3
% type Wiki

本文根据 CC-BY-SA 协议转载翻译自维基百科\href{https://en.wikipedia.org/wiki/George_Kistiakowsky}{相关文章}。

\begin{figure}[ht]
\centering
\includegraphics[width=6cm]{./figures/a36c1f246ada9580.png}
\caption{} \label{fig_QZJS_1}
\end{figure}
乔治·博格达诺维奇·基斯佳科夫斯基（俄语：Георгий Богданович Кистяковский；乌克兰语：Георгій Богданович Кістяківський，罗马化：Heorhii Bohdanovych Kistiakivskyi；1900年12月1日［旧历11月18日］—1982年12月7日）是一位乌克兰裔美国物理化学家，曾任哈佛大学教授，参与“曼哈顿计划”，并在后期担任美国总统德怀特·D·艾森豪威尔的科学顾问。

他出生于当时俄帝国境内的博亚尔卡（Boyarka）\(^\text{[1]}\)，出身于“一个古老的乌克兰哥萨克家族，该家族在俄国革命前是知识精英阶层的一部分”\(^\text{[2]}\)。在俄国内战期间，基斯佳科夫斯基逃离祖国，前往德国，在柏林大学师从马克斯·博登斯坦（Max Bodenstein），获得物理化学博士学位。1926年他移居美国，1930年加入哈佛大学任教，并于1933年成为美国公民。

第二次世界大战期间，基斯佳科夫斯基担任国家国防研究委员会（NDRC）爆炸物开发部门负责人，同时兼任爆炸物研究实验室（ERL）技术主任，负责监督包括RDX与HMX在内的新型炸药的研制。他还参与了爆炸的流体动力学理论研究以及成形装药技术的发展。1943年10月，他被聘为“曼哈顿计划”的顾问，不久后接管X部门，负责研发用于内爆型核武器的爆炸透镜。1945年7月，他亲眼见证了“三位一体”核试验中的首次原子爆炸。数周后，另一枚内爆型核武器“胖子”（Fat Man）被投放至长崎。

1962年至1965年间，基斯佳科夫斯基担任美国国家科学院科学、工程与公共政策委员会（COSEPUP）主席，并于1965年至1973年间出任该院副院长。由于反对越南战争，他切断了与政府的所有联系，并积极投身反战组织“宜居世界委员会”（Council for a Livable World），并于1977年出任该组织主席。
\subsection{早年生活}
乔治·博格达诺维奇·基斯佳科夫斯基于1900年12月1日［旧历11月18日］出生在俄国帝国基辅省的博亚尔卡（今属乌克兰）\(^\text{[1][2][3]}\)。他的祖父亚历山大·费多罗维奇·基斯佳科夫斯基是俄国帝国的法学教授兼律师，专攻刑法\(^\text{[5]}\)。其父博格丹·基斯佳科夫斯基是基辅大学的法理学教授\(^\text{[4]}\)，并于1919年当选为乌克兰国家科学院院士\(^\text{[6]}\)。母亲名为玛丽亚·别连德施塔姆（Maria Berendshtam），他还有一位哥哥亚历山大，后来成为鸟类学家\(^\text{[4]}\)。乔治的叔叔伊戈尔·基斯佳科夫斯基曾任乌克兰国政府内政部长\(^\text{[5]}\)。

基斯佳科夫斯基曾在基辅和莫斯科的私立学校就读，直到1917年俄国革命爆发。随后他加入了反共的白军。1920年，他乘坐一艘被征用的法国船只逃离俄国。此后，他辗转土耳其和南斯拉夫，最终前往德国，并于同年进入柏林大学学习\(^\text{[4]}\)。1925年，他在马克斯·博登斯坦（Max Bodenstein）的指导下获得物理化学博士学位，论文题目为“氯一氧化物与臭氧的光化学分解”。之后，他成为博登斯坦的研究助理\(^\text{[7][4]}\)。他发表的前两篇论文均基于博士论文的延伸研究，并与博登斯坦合作撰写\(^\text{[8]}\)。

1926年，基斯佳科夫斯基以国际教育委员会奖学金获得者的身份赴美。另一位博登斯坦的学生休·斯托特·泰勒（Hugh Stott Taylor）\(^\text{[9]}\)接受了导师对基斯佳科夫斯基的高度评价，为他提供了普林斯顿大学的研究岗位。同年，基斯佳科夫斯基与一位瑞典路德宗女性希尔德加德·梅比乌斯（Hildegard Moebius）结婚\(^\text{[4][10]}\)。1928年，他们的女儿维拉（Vera）出生。1972年，维拉成为麻省理工学院首位获聘物理学教授的女性\(^\text{[11]}\)。1927年，他的两年奖学金到期后，获得了研究助理职位与杜邦奖学金。1928年10月25日，他被任命为普林斯顿大学副教授(^7[4])。在此期间，他与泰勒共同发表了一系列论文\(^\text{[8]}\)，并在泰勒的鼓励下，出版了美国化学学会专著《光化学过程》(^8[4])。

1930年，基斯佳科夫斯基加入哈佛大学任教，并在此度过了整个学术生涯。在哈佛期间，他的研究领域包括热力学、光谱学和化学动力学。同时，他也越来越多地参与政府与工业界的技术咨询。1933年，他在哈佛大学晋升为副教授，并于同年成为美国公民，入选美国艺术与科学院\(^\text{[12]}\)。1938年，他被任命为“阿博特与詹姆斯·劳伦斯化学教授”\(^\text{[13]}\)。翌年，他当选为美国国家科学院院士\(^\text{[14]}\)。1940年，他又当选为美国哲学学会成员\(^\text{[15]}\)。
\subsection{第二次世界大战}
\subsubsection{国家国防研究委员会}
\begin{figure}[ht]
\centering
\includegraphics[width=6cm]{./figures/069a0bbe074c9f6a.png}
\caption{} \label{fig_QZJS_2}
\end{figure}
预见到科学将在尚未参战的美国的二战中发挥更大作用，富兰克林·D·罗斯福总统于1940年6月27日成立了国家国防研究委员会（National Defense Research Committee，NDRC），由万尼瓦·布什（Vannevar Bush）出任主席。哈佛大学校长詹姆斯·B·康南特（James B. Conant）被任命为B部（负责炸弹、燃料、毒气与化学品）的负责人。他任命基斯佳科夫斯基负责A-1小组，A-1小组负责炸药研究。1941年6月，NDRC并入科学研究与发展办公室（Office of Scientific Research and Development，OSRD）。布什成为OSRD主席，康南特接替他出任NDRC主席，基斯佳科夫斯基则成为B部的负责人。1942年12月的一次重组中，B部被拆分，他出任第8部（负责炸药与推进剂）负责人，并一直担任该职直到1944年2月。

基斯佳科夫斯基对美国在炸药与推进剂方面的研究现状感到不满。康南特于1940年10月在宾夕法尼亚州布鲁塞顿（Bruceton）靠近矿务局实验室处设立了爆炸物研究实验室（Explosives Research Laboratory，ERL），基斯佳科夫斯基最初对其活动进行监督并偶有来访；但康南特直到1941年春才正式任命他为该实验室的技术主任。尽管起初受制于设施不足，ERL的规模从1941年的5名工作人员增长到1945年战时的最高峰——162名全职实验室人员。RDX是一项重要的研究方向。这种高能炸药在战前已由德国人开发，挑战在于建立一种可大规模工业生产的工艺。RDX还常与TNT混合制成B型炸药（Composition B），广泛用于各类弹药；以及用于鱼雷与深水炸弹的Torpex。试点工厂在1942年5月投入运行，1943年开始大规模生产。
